\documentclass[hidelinks]{article}
\usepackage{stex}

\libinput{preambleCS}

\title{IsaVODEs: Interactive Verification of Cyber-Physical Systems at Scale}
\author{}
\date{}

\begin{document}

\maketitle

\begin{smodule}[title={IsaVODEs: Interactive Verification of Cyber-Physical Systems at Scale}]{IsabelleEquations}

\symdecl{IsaVODEs}[name=IsaVODEs]
\symdecl{CyberPhysicalSystem}[name={cyber-physical system}]
\symdecl{HybridSystem}[name={hybrid system}]
\symdecl{ODE}[name={ordinary differential equation}]
\symdecl{HybridProgram}[name={hybrid program}]
\symdecl{DiscreteAssignment}[name={discrete assignment}]
\symdecl{ContinuousEvolution}[name={continuous evolution}]
\symdecl{Test}[name=test]
\symdecl{SequentialComposition}[name={sequential composition}]
\symdecl{Choice}[name=choice]
\symdecl{Repeat}[name={finite iteration}]
\symdecl{InteractiveTheoremProver}[name={interactive theorem prover}]
\symdecl{ShallowEmbedding}[name={shallow embedding}]
\symdecl{Flow}[name=flow]
\symdecl{VectorField}[name={vector field}]
\symdecl{InitialValueProblem}[name={initial value problem}]
\symdecl{PicardLindelofTheorem}[name={Picard--Lindelof theorem}]
\symdecl{StateTransformer}[name={state transformer}]
\symdecl{PredicateTransformer}[name={predicate transformer}]
\symdecl{ForwardBox}[name={forward box}]
\symdecl{WeakestLiberalPrecondition}[name={weakest liberal precondition}]
\symdecl{HoareTriple}[name={Hoare triple}]
\symdecl{ForwardDiamond}[name={forward diamond}]
\symdecl{HybridStore}[name={hybrid store}]
\symdecl{Lens}[name=lens]
\symdecl{Dataspace}[name=dataspace]
\symdecl{LiftedExpression}[name={lifted expression}]
\symdecl{LensSubstitution}[name={lens substitution}]
\symdecl{Frame}[name=frame]
\symdecl{NonModification}[name=non-modification]
\symdecl{FramedEvolutionCommand}[name={framed evolution command}]
\symdecl{FramedFrechetDerivative}[name={framed Frechet derivative}]
\symdecl{DifferentialInvariant}[name={differential invariant}]
\symdecl{DifferentialInduction}[name={differential induction}]
\symdecl{DifferentialGhost}[name={differential ghost}]
\symdecl{DarbouxRule}[name={Darboux rule}]
\symdecl{VerificationConditionGeneration}[name={verification condition generation}]
\symdecl{LocalFlow}[name={local flow}]
\symdecl{LipschitzContinuity}[name={Lipschitz continuity}]
\symdecl*{Eisbach}
\symdecl*{Sledgehammer}
\symdecl{ComputerAlgebraSystem}[name={computer algebra system}]
\symdecl{WolframEngine}[name={Wolfram Engine}]
\symdecl*{SageMath}
\symdecl{DerivativeTest}[name={derivative test}]
\symdecl{ReachabilityAnalysis}[name={reachability analysis}]
\symdecl{DeductiveVerification}[name={deductive verification}]
\symdecl*{ARCH22}

\section{Defined Vocabulary}

\begin{sdefinition}[for=IsaVODEs]
    \definame{IsaVODEs} is the Isabelle Verification with Ordinary Differential Equations framework for modelling and verifying cyber-physical systems in Isabelle/HOL.
\end{sdefinition}

\begin{sdefinition}[for=CyberPhysicalSystem]
    A \definame{CyberPhysicalSystem} is a system whose software controller interacts with a physical environment, usually combining discrete program steps with continuous dynamics.
\end{sdefinition}

\begin{sdefinition}[for=HybridSystem]
    A \definame{HybridSystem} is a mathematical model that combines discrete transitions and continuous evolution, commonly represented by hybrid programs and systems of ordinary differential equations.
\end{sdefinition}

\begin{sdefinition}[for=ODE]
    An \definame{ODE} specifies continuous dynamics by relating the derivative of a state trajectory to a vector field.
\end{sdefinition}

\begin{sdefinition}[for=HybridProgram]
    A \definame{HybridProgram} is a program model for hybrid systems with discrete assignments, tests, sequencing, choice, iteration, and continuous evolution commands.
\end{sdefinition}

\begin{sdefinition}[for=DiscreteAssignment]
    A \definame{DiscreteAssignment} instantly updates a program variable with the value of an expression.
\end{sdefinition}

\begin{sdefinition}[for=ContinuousEvolution]
    A \definame{ContinuousEvolution} follows an ordinary differential equation for a nondeterministically chosen duration while respecting a guard or evolution-domain constraint.
\end{sdefinition}

\begin{sdefinition}[for=Test]
    A \definame{Test} checks a predicate in the current state and aborts the current execution branch if the predicate is false.
\end{sdefinition}

\begin{sdefinition}[for=SequentialComposition]
    A \definame{SequentialComposition} executes one hybrid program and then continues with another from each resulting state.
\end{sdefinition}

\begin{sdefinition}[for=Choice]
    A \definame{Choice} permits either of two hybrid programs to execute from the current state.
\end{sdefinition}

\begin{sdefinition}[for=Repeat]
    A \definame{Repeat} executes a hybrid program any finite number of times.
\end{sdefinition}

\begin{sdefinition}[for=InteractiveTheoremProver]
    An \definame{InteractiveTheoremProver} is a proof assistant, such as Isabelle/HOL, Coq, Lean, or PVS, whose trusted kernel checks formal proofs while users and automation construct them.
\end{sdefinition}

\begin{sdefinition}[for=ShallowEmbedding]
    A \definame{ShallowEmbedding} represents the syntax and semantics of a modelling language directly with the host logic of the theorem prover, so that the prover's libraries and automation are immediately available.
\end{sdefinition}

\begin{sdefinition}[for=Flow]
    A \definame{Flow} is a continuously differentiable monoid action \(\phi : T \to C \to C\) on a continuous state space \(C\), satisfying
    \[
        \phi(t_1+t_2)=\phi(t_1)\circ \phi(t_2)
        \qquad\text{and}\qquad
        \phi(0)=id_C.
    \]
\end{sdefinition}

\begin{sdefinition}[for=VectorField]
    A \definame{VectorField} is a function \(f : T \to C \to C\) assigning to each point in space-time the velocity vector that determines an ordinary differential equation.
\end{sdefinition}

\begin{sdefinition}[for=InitialValueProblem]
    An \definame{InitialValueProblem} for a vector field \(f\) asks for a solution \(X:T\to C\) with \(X'(t)=f(t)(X(t))\) and \(X(t_0)=c\) for a specified initial time and state.
\end{sdefinition}

\begin{sdefinition}[for=PicardLindelofTheorem]
    The \definame{PicardLindelofTheorem} gives local uniqueness of solutions to initial value problems when the vector field is Lipschitz continuous on the relevant domain.
\end{sdefinition}

\begin{sdefinition}[for=StateTransformer]
    A \definame{StateTransformer} for a hybrid program is a nondeterministic function \(\alpha:S\to\mathcal{P}S\) mapping an initial state to the set of possible final states.
\end{sdefinition}

\begin{sdefinition}[for=PredicateTransformer]
    A \definame{PredicateTransformer} maps postconditions to predicates over initial states, and is used in IsaVODEs to generate verification conditions for hybrid programs.
\end{sdefinition}

\begin{sdefinition}[for=ForwardBox]
    The \definame{ForwardBox} is the dynamic-logic predicate transformer \(|\alpha]Q\), true at a state \(s\) when every terminating execution of \(\alpha\) from \(s\) ends in a state satisfying \(Q\).
\end{sdefinition}

\begin{sdefinition}[for=WeakestLiberalPrecondition]
    A \definame{WeakestLiberalPrecondition} is the same forward-box transformer viewed as a partial-correctness rule: it imposes \(Q\) on all final states if the program terminates.
\end{sdefinition}

\begin{sdefinition}[for=HoareTriple]
    A \definame{HoareTriple} \(\{P\}\alpha\{Q\}\) states partial correctness: every terminating execution of program \(\alpha\) from a state satisfying \(P\) ends in a state satisfying \(Q\).
\end{sdefinition}

\begin{sdefinition}[for=ForwardDiamond]
    The \definame{ForwardDiamond} is the reachability predicate transformer \(|\alpha\rangle Q\), true at a state \(s\) when some execution of \(\alpha\) from \(s\) reaches a state satisfying \(Q\).
\end{sdefinition}

\begin{sdefinition}[for=HybridStore]
    A \definame{HybridStore} is the state model used by IsaVODEs, allowing continuous variables such as reals, vectors, and matrices to coexist with discrete data structures.
\end{sdefinition}

\begin{sdefinition}[for=Lens]
    A \definame{Lens} \(x::V\Rightarrow S\) consists of get and put functions for accessing and updating a component of a store, subject to the standard lens laws.
\end{sdefinition}

\begin{sdefinition}[for=Dataspace]
    A \definame{Dataspace} is the IsaVODEs declaration mechanism that creates a local hybrid-store context with constants, assumptions, variables, lenses, and independence facts.
\end{sdefinition}

\begin{sdefinition}[for=LiftedExpression]
    A \definame{LiftedExpression} is a user-level term translated into a semantic expression \(S\to V\), hiding explicit store arguments and lens get operations from proof obligations.
\end{sdefinition}

\begin{sdefinition}[for=LensSubstitution]
    A \definame{LensSubstitution} is a semantic store update such as \([x\leadsto e,y\leadsto f]\), used to specify assignments, flows, and vector fields in the lens-based store.
\end{sdefinition}

\begin{sdefinition}[for=Frame]
    A \definame{Frame} is the set of store variables that a program or evolution command may modify; variables outside the frame remain unchanged.
\end{sdefinition}

\begin{sdefinition}[for=NonModification]
    \Definame{NonModification} is the property \(\alpha\;nmods\;A\), meaning that program \(\alpha\) leaves the variables or expressions described by \(A\) unchanged.
\end{sdefinition}

\begin{sdefinition}[for=FramedEvolutionCommand]
    A \definame{FramedEvolutionCommand} lifts ODE behaviour on the continuous part of a store to the whole hybrid store, leaving variables outside the continuous frame fixed.
\end{sdefinition}

\begin{sdefinition}[for=FramedFrechetDerivative]
    A \definame{FramedFrechetDerivative} differentiates an expression with respect to a chosen continuous frame and vector field, ignoring unrelated and discrete store components.
\end{sdefinition}

\begin{sdefinition}[for=DifferentialInvariant]
    A \definame{DifferentialInvariant} is a predicate preserved by an evolution command, equivalently an invariant set for the associated dynamical system.
\end{sdefinition}

\begin{sdefinition}[for=DifferentialInduction]
    \Definame{DifferentialInduction} proves differential invariants by comparing the derivatives of expressions occurring in equalities and inequalities.
\end{sdefinition}

\begin{sdefinition}[for=DifferentialGhost]
    A \definame{DifferentialGhost} is an auxiliary differential equation added to a system to make an invariant proof possible while preserving the original system's correctness claim.
\end{sdefinition}

\begin{sdefinition}[for=DarbouxRule]
    A \definame{DarbouxRule} is a specialised consequence of the differential ghost rule for proving invariance of equalities or inequalities whose derivative has a Darboux-style factor.
\end{sdefinition}

\begin{sdefinition}[for=VerificationConditionGeneration]
    \Definame{VerificationConditionGeneration} reduces a program-correctness claim to logical and mathematical proof obligations that can be solved by simplification, tactics, or external proof search.
\end{sdefinition}

\begin{sdefinition}[for=LocalFlow]
    A \definame{LocalFlow} is the Isabelle/HOL predicate used by IsaVODEs to certify that a supplied flow solves a framed ODE on a given domain.
\end{sdefinition}

\begin{sdefinition}[for=LipschitzContinuity]
    \definame{LipschitzContinuity} is the regularity condition used to invoke uniqueness of ODE solutions, and is automated in IsaVODEs through \(C^1\)-based certification.
\end{sdefinition}

\begin{sdefinition}[for=Eisbach]
    \definame{Eisbach} is Isabelle's proof-method language, used in IsaVODEs to package high-level tactics such as \texttt{wlp\_solve}, \texttt{dInv}, and \texttt{dInduct\_mega}.
\end{sdefinition}

\begin{sdefinition}[for=Sledgehammer]
    \definame{Sledgehammer} is Isabelle's bridge to external automated theorem provers, used to search for proofs that are reconstructed and checked inside Isabelle.
\end{sdefinition}

\begin{sdefinition}[for=ComputerAlgebraSystem]
    A \definame{ComputerAlgebraSystem} is an external symbolic algebra tool used as an oracle for candidate ODE solutions, which IsaVODEs then attempts to certify in Isabelle.
\end{sdefinition}

\begin{sdefinition}[for=WolframEngine]
    The \definame{WolframEngine} is the Mathematica-family computer algebra backend used by IsaVODEs to suggest symbolic solutions through Wolfram expressions and \texttt{DSolve}.
\end{sdefinition}

\begin{sdefinition}[for=SageMath]
    \definame{SageMath} is the open-source computer algebra backend integrated by IsaVODEs through a Python-facing API and an intermediate ODE representation.
\end{sdefinition}

\begin{sdefinition}[for=DerivativeTest]
    A \definame{DerivativeTest} is a calculus theorem for detecting monotonicity or local extrema from first and second derivatives, formalised in IsaVODEs to support arithmetic side proofs.
\end{sdefinition}

\begin{sdefinition}[for=ReachabilityAnalysis]
    \Definame{ReachabilityAnalysis} verifies hybrid systems by over-approximating or computing reachable state sets until a violation or fixed point is found.
\end{sdefinition}

\begin{sdefinition}[for=DeductiveVerification]
    \Definame{DeductiveVerification} establishes correctness by formal proof from semantic rules, as in IsaVODEs, KeYmaera X, HHL, and related theorem-proving frameworks.
\end{sdefinition}

\begin{sdefinition}[for=ARCH22]
    \definame{ARCH22} refers to the 2022 Applied Verification of Continuous and Hybrid Systems friendly competition benchmark set used to evaluate IsaVODEs.
\end{sdefinition}

\section{Scope and Contributions}

Huerta y Munive, Foster, Gleirscher, Struth, Pardillo Laursen, and Hickman present \sn{IsaVODEs} as an Isabelle/HOL framework for verifying \sns{CyberPhysicalSystem}.
The starting point is the standard modelling split in which a controller is a discrete, state-updating, and possibly nondeterministic program, while the physical environment is specified by systems of \sns{ODE}.
The paper develops this into a scalable \sn{DeductiveVerification} framework for \sns{HybridSystem}, using \sns{HybridProgram}, predicate-transformer semantics, and Isabelle's libraries for analysis and ODEs.

The framework is a \sn{ShallowEmbedding} in Isabelle/HOL.
This choice makes the language extensible and allows proofs to use Isabelle's typed higher-order logic, existing mathematical libraries, \sn{Sledgehammer}, and ordinary Isabelle proof methods.
The paper positions this against domain-specific deductive tools such as KeYmaera X: IsaVODEs seeks comparable automation while preserving the trust model and openness of a general \sn{InteractiveTheoremProver}.

The main contributions are a modelling language for hybrid programs, a lens-based \sn{HybridStore}, frame and local-reasoning laws, \sn{FramedFrechetDerivative}s for ODE reasoning over part of the store, proof methods for \sn{VerificationConditionGeneration}, two ODE workflows based on flows or invariants, integration with \sns{ComputerAlgebraSystem}, and an \sn{ARCH22} benchmark evaluation.
The paper also adds forward diamonds for progress and reachability claims, generalised \sns{DifferentialGhost} and \sns{DarbouxRule}, automated \sn{LipschitzContinuity} certification, and formalised \sns{DerivativeTest}.

\section{Hybrid-System Semantics}

\subsection{Dynamical Systems}

The mathematical basis distinguishes explicit and implicit descriptions of continuous dynamics.
An explicit description uses a \sn{Flow} \(\phi:T\to C\to C\), where \(T\) is a nondiscrete submonoid of \(\mathbb{R}\) and \(C\) is a continuous state space.
For an initial state \(c\in C\), the trajectory \(\phi_c:T\to C\) is given by \(\phi_c(t)=\phi(t)(c)\).
The orbit map sends an initial state to all states along its trajectory:
\[
    \gamma_\phi(c)=\{\phi_c(t)\mid t\in T\}.
\]

An implicit description uses a \sn{VectorField} \(f:T\to C\to C\).
A solution to the associated system of \sns{ODE} is a \(C^1\) function \(X:T\to C\) satisfying
\[
    X'(t)=f(t)(X(t)).
\]
An \sn{InitialValueProblem} additionally fixes \(X(t_0)=c\).
The \sn{PicardLindelofTheorem} is the uniqueness principle used later: under \sn{LipschitzContinuity}, the solutions of an initial value problem coincide locally, and when this holds uniformly the solutions induce a \sn{Flow}.

\subsection{State Transformers}

The semantic core of \sn{IsaVODEs} represents \sns{HybridProgram} as nondeterministic \sns{StateTransformer} \(\alpha:S\to\mathcal{P}S\).
This semantic view mirrors the usual grammar
\[
    \alpha ::= x := e \mid x' = f\ \&\ G \mid ?P \mid \alpha;\alpha
        \mid \alpha\cup\alpha \mid \alpha^* .
\]
Here \(x := e\) is a \sn{DiscreteAssignment}, \(x'=f\ \&\ G\) is a \sn{ContinuousEvolution} with guard \(G\), \(?P\) is a \sn{Test}, \(\alpha;\beta\) is \sn{SequentialComposition}, \(\alpha\cup\beta\) is nondeterministic \sn{Choice}, and \(\alpha^*\) is \sn{Repeat}.

The state-transformer equations are the expected Kleisli operations for the powerset monad:
\[
    (\alpha;\beta)(s)=\bigcup_{s'\in\alpha(s)}\beta(s'),
    \qquad
    (\alpha\cup\beta)(s)=\alpha(s)\cup\beta(s),
\]
\[
    \alpha^*(s)=\bigcup_{i\in\mathbb{N}}\alpha^i(s),
    \qquad
    \alpha^0=skip,\quad \alpha^{i+1}=\alpha;\alpha^i .
\]
Assertions are identified with subidentity state transformers, so predicates, sets of states, and tests can be moved between program and logical viewpoints.

\subsection{Stores and Lenses}

\sn{IsaVODEs} models program variables with \sns{Lens}.
A lens \(x::V\Rightarrow S\) has \(get_x:S\to V\) and \(put_x:V\to S\to S\), satisfying
\[
    get_x(put_x(v)(s))=v,\qquad
    put_x(v)\circ put_x(v')=put_x(v),\qquad
    put_x(get_x(s))(s)=s.
\]
The source \(S\) is the \sn{HybridStore}; the view \(V\) is the component selected by the lens.
An expression is a function \(e:S\to V\), and an assignment is represented by
\[
    (x:=e)=\langle \lambda s.\ put_x(e(s))(s)\rangle,
\]
where angle brackets indicate the deterministic function lifted into a nondeterministic state transformer.

For continuous dynamics, a guarded orbit map generalises the orbit of a \sn{Flow} to all solutions of a \sn{VectorField}.
For a guard \(G\), interval assignment \(U\), initial time \(t_0\), and state \(c\), the paper uses the set of all \(X(t)\) such that \(X\) solves the \sn{InitialValueProblem}, \(t\in U(c)\), and the trajectory remains inside \(G\) up to \(t\).
This supplies the semantic content of \sns{ContinuousEvolution}; later sections lift it from the continuous space to the whole \sn{HybridStore} by framing.

\section{Predicate Transformers}

\subsection{Forward Boxes and Hoare Triples}

The \sn{ForwardBox}, also called the \sn{WeakestLiberalPrecondition}, is the \sn{PredicateTransformer}
\[
    |\alpha]Q(s) \Longleftrightarrow
    \forall s'.\, s'\in\alpha(s)\Longrightarrow Q(s').
\]
It supports partial-correctness reasoning: if \(\alpha\) terminates from \(s\), all resulting states satisfy \(Q\).
The standard equations become rewrite rules for \sn{VerificationConditionGeneration}:
\[
\begin{array}{rclcrcl}
|skip]Q &=& Q, &&
|abort]Q &=& \top,\\
|?P]Q &=& P\Longrightarrow Q, &&
|x:=e]Q &=& Q[e/x],\\
|\alpha;\beta]Q &=& |\alpha]|\beta]Q, &&
|\alpha\cup\beta]Q &=& |\alpha]Q\wedge|\beta]Q,\\
|if\ T\ then\ \alpha\ else\ \beta]Q
&=& (T\Longrightarrow|\alpha]Q)\wedge(\neg T\Longrightarrow|\beta]Q).
\end{array}
\]
Loops require invariant rules rather than simple simplification.

A \sn{HoareTriple} is defined from the \sn{ForwardBox}:
\[
    \{P\}\alpha\{Q\}\Longleftrightarrow P\Longrightarrow|\alpha]Q.
\]
The resulting Hoare rules include rules for skip, abort, tests, assignments, sequential composition, choice, conditionals, loops, and evolution commands.
For ODEs, the important split is methodological: when a certified \sn{Flow} is available, use the flow law; when it is not available, provide a \sn{DifferentialInvariant}.

\subsection{Forward Diamonds}

The \sn{ForwardDiamond} extends the framework from safety-style partial correctness to reachability and progress:
\[
    |\alpha\rangle Q(s) \Longleftrightarrow
    \exists s'.\, s'\in\alpha(s)\wedge Q(s').
\]
It satisfies the duality \( |\alpha\rangle Q=\neg|\alpha]\neg Q\), and therefore has laws corresponding to the forward-box laws:
\[
\begin{array}{rclcrcl}
|skip\rangle Q &=& Q, &&
|abort\rangle Q &=& \bot,\\
|?P\rangle Q &=& P\wedge Q, &&
|x:=e\rangle Q &=& Q[e/x],\\
|\alpha;\beta\rangle Q &=& |\alpha\rangle|\beta\rangle Q, &&
|\alpha\cup\beta\rangle Q &=& |\alpha\rangle Q\vee|\beta\rangle Q.
\end{array}
\]
For \sns{ContinuousEvolution}, the diamond law existentially quantifies a solution and a time at which the postcondition is reached while the guard has held along the preceding trajectory.

\section{Hybrid Modelling Language}

\sn{Dataspace}s are IsaVODEs' user-facing store declarations.
A declaration introduces constants, named assumptions, and variables, and internally creates an Isabelle locale equipped with lenses and independence assumptions.
Because variables are \sns{Lens}, the resulting model supports local reasoning, inheritance of store contexts, and typed variables over reals, vectors, matrices, and discrete data types.

\sn{LiftedExpression}s bridge the gap between ordinary mathematical syntax and semantic functions \(S\to V\).
The user writes expressions such as \(x^2+y^2\leq c^2\), while Isabelle stores a function involving \(get_x\), \(get_y\), and a hidden state argument.
The translation marks lifted terms with expression syntax, inserts getters for lenses, applies existing expressions to the state, and leaves constants unchanged.
The pretty-printer reverses this translation, so proof goals remain readable.

\sn{LensSubstitution}s provide the notation used for assignments, flows, and framed vector fields:
\[
    [x_1\leadsto e_1,x_2\leadsto e_2,\ldots]
    = id(x_1\leadsto e_1)(x_2\leadsto e_2)\cdots .
\]
Variables not mentioned in the substitution are unchanged.
Substitution into expressions is written \(e[v/x]=e\circ[x\leadsto v]\), giving the familiar \(Q[e/x]\) law for assignments.
An Isabelle simplification procedure reorders and reduces substitutions during verification-condition generation.

Vectors and matrices are inherited from HOL-Analysis finite Cartesian products.
IsaVODEs adds surface syntax for matrices and component access through composed \sns{Lens}.
If \(p::\mathbb{R}^n\Rightarrow S\), then coordinate lenses are obtained by composing \(p\) with the \(i\)-th vector-coordinate lens.
This lets a specification speak both about structured engineering objects and about their individual coordinates.

\section{Local Reasoning}

\subsection{Frames and Non-modification}

A \sn{Frame} is represented by lens sums.
Independent lenses \(x_1::V_1\Rightarrow S\) and \(x_2::V_2\Rightarrow S\) combine as
\[
    x_1\oplus x_2::V_1\times V_2\Rightarrow S,
\]
with the view obtained by pairing the two getters and the update obtained by composing the two puts.
This represents finite heterogeneous variable sets, up to the usual product isomorphisms.

The paper uses lens preorders and equivalences to reason about inclusion and equality of variable sets.
It also uses unrestriction, written \(A\sharp e\), to express that an expression \(e\) does not semantically depend on variables in \(A\).
\sn{NonModification}, written \(\alpha\;nmods\;A\), says that all executions of \(\alpha\) preserve the variables or expressions described by \(A\).
Tests modify nothing, assignments modify only their assigned variables, and program combinators inherit non-modification from their parts.

The central frame rules are:
\[
    \frac{\{P\}\alpha\{Q\}\qquad \alpha\;nmods\;I}
         {\{P\wedge I\}\alpha\{Q\wedge I\}},
\]
and the variant that infers immutability from a frame \(A\) plus the condition that \(I\) depends only on variables outside \(A\).
These rules support separation-logic-style local reasoning over hybrid stores.

\subsection{Framed Evolution Commands}

\sn{FramedEvolutionCommand}s turn continuous dynamics over a part \(C\) of a store into state transformers over the full store \(S\).
Given a frame lens \(x::C\Rightarrow S\), a predicate \(G:S\to\mathbb{B}\) restricts to \(G^s_x:C\to\mathbb{B}\) by
\[
    G^s_x(c) \Longleftrightarrow G(put_x(c)(s)).
\]
A state-wide substitution \(f:T\to S\to S\) frames to a vector field on \(C\) by
\[
    f^s_x(t)(c)=get_x(f(t)(put_x(c)(s))).
\]
The evolution command then updates only the continuous frame \(x\), by placing each solution point \(X(t)\) back into the original store with \(put_x(X(t))(s)\).
Variables outside the frame remain unchanged, which gives ODE non-modification laws for free.

\subsection{Differential Invariants and Framed Derivatives}

The predicate \(diff\_inv\ x\ f\ G\ U\ t_0\ I\) abbreviates that \(I\) is a \sn{DifferentialInvariant} for a framed evolution command.
In Hoare form, it states
\[
    diff\_inv\ x\ f\ G\ U\ t_0\ I
    \Longleftrightarrow
    \{I\}\ (x'=f\ \&\ G)^ {t_0}_U\ \{I\}.
\]
\sn{DifferentialInduction} reduces invariance of equalities and inequalities to derivative comparisons.

For an expression \(e:S\to U\), vector field \(f:S\to S\), and frame \(x::C\Rightarrow S\), the \sn{FramedFrechetDerivative} is
\[
    (D^f_x e)(s)
    =
    D(e^s_x)(get_x(s))(get_x(f(s))).
\]
It is the directional Frechet derivative of the framed expression in the direction induced by the ODE.
The core invariant rules are:
\[
\begin{array}{rcl}
G\Longrightarrow D^f_x e_1=D^f_x e_2
&\Longrightarrow&
diff\_inv\ x\ f\ G\ \mathbb{R}_{\geq 0}\ 0\ (e_1=e_2),\\
G\Longrightarrow D^f_x e_1\leq D^f_x e_2
&\Longrightarrow&
diff\_inv\ x\ f\ G\ \mathbb{R}_{\geq 0}\ 0\ (e_1\leq e_2),\\
G\Longrightarrow D^f_x e_1\leq D^f_x e_2
&\Longrightarrow&
diff\_inv\ x\ f\ G\ \mathbb{R}_{\geq 0}\ 0\ (e_1<e_2).
\end{array}
\]
The accompanying derivative laws say that constants and variables outside the frame have derivative \(0\), variables inside the frame use the ODE right-hand side, and sums, products, powers, and logarithms follow the usual calculus rules.

\subsection{Ghosts and Darboux Rules}

\sn{DifferentialInduction} is incomplete for many invariants.
\sn{IsaVODEs} therefore formalises a framed \sn{DifferentialGhost} rule, adding fresh auxiliary ODEs to prove a property of the original system.
In schematic form:
\[
    \frac{\{P\}\{x'=f,\ y'=A\cdot y+b\ \&\ G\}\{Q\}}
         {\{\exists v.\ P[v/y]\}\{x'=f\ \&\ G\}\{\exists v.\ Q[v/y]\}} .
\]
The formal rule requires independence between the ghost variables and the original frame.
The paper then derives \sns{DarbouxRule} for equality and inequality invariants, giving specialised ghost-based rules for cases where the derivative of an expression can be related to the expression itself.

\section{Reasoning Components}

\sn{IsaVODEs} packages the previous semantic laws into \sn{Eisbach} proof methods.
The method \texttt{vderiv} recursively certifies vector derivative equations by applying derivative rules and higher-order unification.
The paper extends this tactic with rules for exponentials, square roots, tangent, cotangent, vector inner products, and norms.

\sn{LipschitzContinuity} certification is automated through the formal theorem that \(C^1\) functions are locally Lipschitz continuous:
\[
    open(S)\wedge open(T)\wedge
    \bigl(\forall \tau\in T.\forall s\in S.\ D(f(\tau))\text{ is the derivative at }s\bigr)
    \wedge continuous\_on(S,D)
    \Longrightarrow local\_lipschitz(T,S,f).
\]
The tactics \texttt{c1\_lipschitz} and \texttt{c1\_lipschitzI} apply this theorem and attempt to discharge openness, continuity, and differentiability obligations.
The method \texttt{local\_flow\_on\_auto} combines \texttt{c1\_lipschitz} and \texttt{vderiv} to certify that a supplied flow is a \sn{LocalFlow}.

The main flow-based verification methods are:
\[
\begin{array}{ll}
\texttt{wlp\_simp} & \text{simplifies forward-box goals using wlp laws and loop invariants},\\
\texttt{wlp\_flow} & \text{uses a previously proved local-flow theorem},\\
\texttt{wlp\_solve} & \text{uses a candidate flow and calls local-flow certification},\\
\texttt{wlp\_full} & \text{tries to close the residual real-arithmetic goals}.
\end{array}
\]
Additional simplification tactics support algebraic rearrangement of monomials and powers.

The \sn{ComputerAlgebraSystem} integration supplies candidate solutions, never unchecked final proofs.
The command \texttt{find\_local\_flow} extracts an ODE from the current goal, translates it into an intermediate arithmetic-expression representation, asks a backend for a symbolic solution, translates the result back into Isabelle syntax, and suggests a \texttt{wlp\_solve} invocation.
\sn{WolframEngine} integration targets Wolfram expressions and \texttt{DSolve}; \sn{SageMath} integration uses a Python-facing API and adds preprocessing for independent equations and higher-order ODE encodings.
Candidate solutions can fail certification if they use functions absent from Isabelle or are undefined on parts of the requested domain.

For invariant-based reasoning, \texttt{dWeaken} applies differential weakening, \texttt{dInduct} applies differential induction and framed derivative simplification, and \texttt{dInduct\_auto} adds expression automation.
The search method \texttt{dInduct\_mega} repeatedly tries known facts, differential weakening, differential cuts, conditional splitting, and automatic differential induction.
\texttt{dInv} wraps this into the common task of proving a supplied invariant sufficient for a Hoare-style ODE goal.
\texttt{dGhost} automates the framed \sn{DifferentialGhost} rule.

The paper also formalises \sns{DerivativeTest} to support real-arithmetic obligations that remain after VCG.
It defines monotonicity and local extrema, proves interval-composition lemmas for increasing and decreasing functions, and proves first- and second-derivative tests.
The first derivative test derives monotonicity from the sign of the derivative on an interval.
The second derivative test derives a local maximum or minimum from a zero first derivative and a negative or positive continuous second derivative.

\section{Evaluation}

The evaluation uses \(66\) \sn{ARCH22} hybrid-systems theorem-proving problems.
The set contains \(5\) regular-program problems, \(32\) continuous-dynamics problems, and \(29\) hybrid-program problems.
The authors classify them by purpose, competition category, and dynamics; \(40\) have linear dynamics, \(21\) have nonlinear dynamics, and \(5\) have no ODEs.

\sn{IsaVODEs} proves \(58\) of the \(66\) benchmark problems fully inside Isabelle.
It proves \(4\) further problems when external arithmetic certification is allowed.
Within the Design Shapes category, the paper reports \(59\) solved out of \(61\), with \(55\) fully proved without a computer algebra oracle.
This is presented as comparable with KeYmaera X and HHLPy on the same competition material.

The authors first attempted flow-based proofs with supplied solutions, solving \(39\) problems that way.
The remaining solved problems required higher-order Isabelle reasoning or dL-style methods such as \sn{DifferentialInduction}, \sns{DifferentialGhost}, differential cuts, weakening, and \sns{DarbouxRule}.
The tactics substantially shorten scripts: \(48\) benchmark problems are solved by a single tactic call.
The largest remaining limitation is real arithmetic, especially nonlinear arithmetic obligations that are too hard for standard SMT-backed automation.

\section{Related Work and Outlook}

The paper distinguishes \sn{ReachabilityAnalysis} from \sn{DeductiveVerification}.
Reachability tools approximate or compute reachable state sets, while deductive tools prove correctness from logical rules.
IsaVODEs belongs to the deductive line and is compared with PVS and Coq formalisations of hybrid systems, Isabelle/HOL verification frameworks, Hybrid Hoare Logic and HHLPy, KeYmaera X and differential dynamic logic, and related formalisation work in Isabelle's Archive of Formal Proofs.

Compared with domain-specific tools, the distinctive IsaVODEs choice is to keep the framework inside Isabelle/HOL.
External tools may suggest facts or solutions, but certification is intended to happen in Isabelle.
The framework benefits from Isabelle's libraries, document model, frontend, syntax translations, code generation support, hammers, and extensible proof-method language.

The conclusion identifies four directions.
First, real-arithmetic decision procedures need better certified integration.
Second, the modal side of the framework can be extended beyond the existing forward diamond and backward diamond work.
Third, the differential ghost rule can be generalised further to match the full flexibility of dL's syntactic rule.
Fourth, IsaVODEs can be connected to software-engineering workflows and engineering notations because its syntax and semantics are extensible within Isabelle.

\end{smodule}
\end{document}
